\documentclass{article}
\usepackage{arxiv}
% ── Paquetes ──────────────────────────────────────────────────────────────────
\usepackage[utf8]{inputenc}
\usepackage[T1]{fontenc}
\usepackage[spanish]{babel}
\usepackage{amsmath, amssymb}
\usepackage{graphicx}
\usepackage{float}
\usepackage{booktabs}
\usepackage{hyperref}
\usepackage{geometry}
\usepackage{listings}
\usepackage{xcolor}
\usepackage{caption}
\usepackage{subcaption}
\usepackage{url}
\usepackage{amsfonts}
\usepackage{nicefrac}
\usepackage{microtype}
\geometry{a4paper, margin=2.5cm}

% ── Estilo de código Python ───────────────────────────────────────────────────
\lstset{
  language=Python,
  basicstyle=\ttfamily\small,
  keywordstyle=\color{blue}\bfseries,
  commentstyle=\color{gray},
  stringstyle=\color{orange},
  numbers=left,
  numberstyle=\tiny\color{gray},
  breaklines=true,
  frame=single,
  backgroundcolor=\color{gray!8}
}
% ─────────────────────────────────────────────────────────────────────────────
\title{TP 2.1 --- GENERADORES PSEUDOALEATORIOS}

\author{
  Santiago Carbó \\
  Universidad Tecnológica Nacional | FRRO\\
  Zeballos 1341, Rosario, Santa Fe  \\
  \texttt{santiago.carbo05@gmail.com} \\
  \And
  Luca Gritti \\
  Universidad Tecnológica Nacional | FRRO\\
  Zeballos 1341, Rosario, Santa Fe \\
  \texttt{lucagritti24@gmail.com} \\
  \And
  Valentín Cotorruelo \\
  Universidad Tecnológica Nacional | FRRO\\
  Zeballos 1341, Rosario, Santa Fe  \\
  \texttt{valentincotorruelo@gmail.com} \\
  \And
  Felipe Sosa Bianciotto \\
  Universidad Tecnológica Nacional | FRRO\\
  Zeballos 1341, Rosario, Santa Fe  \\
  \texttt{felipesosabianciotto@gmail.com} \\
  \And
  Felipe Muntaabski \\
  Universidad Tecnológica Nacional | FRRO\\
  Zeballos 1341, Rosario, Santa Fe  \\
  \texttt{felipemuntaabski@gmail.com} \\
}

% ─────────────────────────────────────────────────────────────────────────────
\begin{document}
\maketitle

% ── Abstract ──────────────────────────────────────────────────────────────────
\begin{abstract}
El presente trabajo implementa y evalúa tres generadores de números
pseudoaleatorios en Python\,3: el método de los \textbf{Cuadrados Medios}
(Von Neumann), el \textbf{Generador Congruencial Lineal} (GCL, parámetros
glibc/POSIX) y el \textbf{Mersenne Twister} incluido en la biblioteca
estándar de Python.  Sobre cada generador se aplica una batería de cinco
pruebas estadísticas — Chi-cuadrado, Kolmogorov-Smirnov, Rachas,
Autocorrelación y Póker — para evaluar su uniformidad e independencia.
Los resultados muestran que el GCL y el Mersenne Twister superan todos los
tests, mientras que los Cuadrados Medios fallan el test del Póker,
evidenciando sus conocidas limitaciones en la distribución de dígitos.
Se presentan gráficos comparativos de histogramas, ECDF, diagramas de
dispersión serial y función de autocorrelación.

\medskip
\noindent\textbf{Palabras clave:}
Simulación $\cdot$ Números pseudoaleatorios $\cdot$ GCL $\cdot$
Cuadrados Medios $\cdot$ Mersenne Twister $\cdot$ Test Chi-cuadrado $\cdot$
Kolmogorov-Smirnov $\cdot$ Test de Rachas $\cdot$ Autocorrelación
\end{abstract}

% ── 1. Introducción ───────────────────────────────────────────────────────────
\section{Introducción}

Un \emph{número pseudoaleatorio} es aquel generado por un proceso
completamente determinista que, no obstante, produce secuencias que no
exhiben ningún patrón o regularidad aparente desde un punto de vista
estadístico~\cite{knuth1997}.  Ante las mismas condiciones iniciales
(semilla), el algoritmo siempre reproduce exactamente la misma secuencia.
Esta característica diferencia a los generadores \emph{pseudoaleatorios}
de los generadores \emph{verdaderamente aleatorios}, que explotan fenómenos
físicos impredecibles (ruido térmico, decaimiento radiactivo, etc.).

La calidad de un generador pseudoaleatorio es determinante en la validez de
cualquier simulación de Monte Carlo, ya que errores sistemáticos en la
generación se propagan directamente a los estimadores obtenidos.  Por ello
resulta imprescindible aplicar pruebas estadísticas formales que permitan
aceptar o rechazar la hipótesis de que los números producidos se comportan
como una muestra de una distribución uniforme $U[0,1)$ con observaciones
independientes.

Este trabajo implementa tres generadores representativos de distintas épocas
y calidades:
\begin{itemize}
  \item \textbf{Cuadrados Medios} (Von Neumann, 1946): generador histórico,
        de período típicamente corto.
  \item \textbf{GCL} (Generador Congruencial Lineal, parámetros POSIX):
        estándar en muchos sistemas durante décadas.
  \item \textbf{Mersenne Twister} (MT19937, Python \texttt{random}):
        referente moderno con período $2^{19937}-1$.
\end{itemize}

% ── 2. Generadores implementados ─────────────────────────────────────────────
\section{Generadores de números pseudoaleatorios}

\subsection{Método de los Cuadrados Medios}

Propuesto por John von Neumann en 1946~\cite{metropolis1987}, es
históricamente el primer generador pseudoaleatorio documentado.  A partir
de una semilla $X_0$ de $d$ dígitos decimales, la recurrencia es:

\begin{equation}
  X_{n+1} = \text{dígitos centrales}\!\left(X_n^2,\; d\right)
  \label{eq:cuadrados}
\end{equation}

La normalización a $[0,1)$ se obtiene como $U_n = X_n / 10^d$.
El método presenta serias limitaciones: el período es generalmente muy
corto y puede degenerar en ciclos triviales (por ejemplo, $X_n = 0$
para todo $n$ si la semilla eleva al cuadrado dando cero en los dígitos
centrales).  Se implementó con $d=8$ dígitos para maximizar su período
dentro de las restricciones del método.

\subsection{Generador Congruencial Lineal (GCL)}

Introducido por D.\,H. Lehmer en 1949~\cite{lehmer1951}, el GCL es el
generador pseudoaleatorio más ampliamente estudiado.  Su recurrencia es:

\begin{equation}
  X_{n+1} = (a \cdot X_n + c) \bmod m
  \label{eq:gcl}
\end{equation}

con normalización $U_n = X_n / m \in [0,1)$.
La elección de parámetros es crítica: según el \emph{Teorema de
Hull-Dobell}~\cite{hull1962}, el período es máximo ($= m$) si y sólo si:
\begin{enumerate}
  \item $\gcd(c, m) = 1$,
  \item $a \equiv 1 \pmod{p}$ para todo primo $p$ que divide a $m$,
  \item si $4 \mid m$, entonces $a \equiv 1 \pmod{4}$.
\end{enumerate}

En este trabajo se utilizan los parámetros del estándar glibc/POSIX,
que garantizan período completo:

\begin{equation}
  m = 2^{31},\quad a = 1\,103\,515\,245,\quad c = 12\,345
  \label{eq:gcl_params}
\end{equation}

\subsection{Mersenne Twister (Python \texttt{random})}

El Mersenne Twister (MT19937)~\cite{matsumoto1998} es el generador
estándar de Python desde la versión~2.3.  Tiene las siguientes
propiedades destacadas:

\begin{itemize}
  \item Período: $2^{19937} - 1 \approx 4.3 \times 10^{6001}$.
  \item Equidistribución en 623 dimensiones.
  \item Supera la mayoría de los tests estadísticos estándar.
\end{itemize}

Dado que su implementación interna no es parte del alcance de este trabajo,
se usa como referencia de calidad mediante la función
\texttt{random.random()} de la biblioteca estándar.

% ── 3. Tests estadísticos ─────────────────────────────────────────────────────
\section{Tests estadísticos}

Se aplicaron cinco tests para evaluar \emph{uniformidad} e
\emph{independencia} de cada generador, con un nivel de significación
$\alpha = 0{,}05$ en todos los casos.  La hipótesis nula $H_0$ en cada test
es que la muestra proviene de una distribución $U[0,1)$ con observaciones
independientes.

\subsection{Test de uniformidad Chi-cuadrado}

Divide el intervalo $[0,1)$ en $k=10$ subintervalos iguales y compara las
frecuencias observadas $O_i$ con la frecuencia esperada $E_i = n/k$:

\begin{equation}
  \chi^2 = \sum_{i=1}^{k} \frac{(O_i - E_i)^2}{E_i}
  \label{eq:chi2}
\end{equation}

Bajo $H_0$, $\chi^2 \sim \chi^2(k-1)$.  El valor crítico para
$k=10$, $\alpha=0{,}05$ es $\chi^2_{9;\,0{,}05} = 16{,}919$.

\subsection{Test de Kolmogorov-Smirnov}

Mide la máxima discrepancia entre la función de distribución empírica
$F_n(x)$ y la teórica $F(x)=x$:

\begin{equation}
  D = \sup_{x \in [0,1)} \left| F_n(x) - x \right|
  \label{eq:ks}
\end{equation}

El valor crítico aproximado para $\alpha=0{,}05$ es
$D^* = 1{,}36/\sqrt{n}$~\cite{massey1951}.
Para $n=1000$: $D^* \approx 0{,}04301$.

\subsection{Test de Rachas}

Evalúa la independencia de la secuencia contando el número de
\emph{rachas} $R$ alrededor de la mediana muestral.  Bajo $H_0$ la
distribución asintótica es normal:

\begin{equation}
  Z = \frac{R - \mathbb{E}[R]}{\sqrt{\mathrm{Var}[R]}},\qquad
  \mathbb{E}[R] = \frac{2\,n_1\,n_2}{n_1+n_2}+1
  \label{eq:rachas}
\end{equation}

donde $n_1$ y $n_2$ son los conteos por encima y por debajo de la mediana.
Se rechaza $H_0$ si $|Z| > 1{,}96$.

\subsection{Test de Autocorrelación}

Calcula el coeficiente de correlación de Pearson entre $X_i$ y
$X_{i+\ell}$ para un retardo $\ell=1$:

\begin{equation}
  r_\ell = \frac{\displaystyle\sum_{i=1}^{n-\ell}(X_i-\bar{X})(X_{i+\ell}-\bar{X})}{(n-\ell)\,s^2}
  \label{eq:autocorr}
\end{equation}

Bajo $H_0$: $r_\ell \sim \mathcal{N}(0,\,1/(n-\ell))$, por lo que
$Z = r_\ell \sqrt{n-\ell}$ tiene distribución $\mathcal{N}(0,1)$.
Se rechaza $H_0$ si $|Z| > 1{,}96$.

\subsection{Test del Póker}

Agrupa los números en bloques de $d=5$, discretiza cada valor a un dígito
decimal ($\lfloor 10\,U_i \rfloor$) y clasifica el bloque según cuántos
valores distintos contiene, análogamente a las manos del póker.  Las
probabilidades teóricas en base 10 son:

\begin{table}[H]
  \centering
  \caption{Categorías y probabilidades teóricas del Test del Póker ($d=5$, base~10).}
  \label{tab:poker}
  \begin{tabular}{lcc}
    \toprule
    Categoría & Valores distintos & $p_{\text{teórica}}$ \\
    \midrule
    5 distintos & 5 & 0{,}30240 \\
    Un par      & 4 & 0{,}50400 \\
    Dos pares   & 3 & 0{,}10800 \\
    Trío        & 3 & 0{,}07200 \\
    Full        & 2 & 0{,}00900 \\
    Póker       & 2 & 0{,}00450 \\
    Quintilla   & 1 & 0{,}00010 \\
    \bottomrule
  \end{tabular}
\end{table}

Las frecuencias observadas se comparan con las esperadas mediante un
estadístico $\chi^2$ con $6$ grados de libertad; el valor crítico para
$\alpha=0{,}05$ es $\chi^2_{6;\,0{,}05} = 12{,}592$.

% ── 4. Resultados ─────────────────────────────────────────────────────────────
\section{Resultados}

Se generaron $n=1000$ números con cada generador usando semilla
$X_0=12\,345$.  Los resultados numéricos de todos los tests se resumen en
la Tabla~\ref{tab:resultados}.

\begin{table}[H]
  \centering
  \caption{Resultados de los tests estadísticos ($n=1000$, $\alpha=0{,}05$).
           ``Estadístico'' es el valor calculado; ``Crítico'' es el umbral
           de rechazo.  \textbf{PASA}: no se rechaza $H_0$.  \textbf{FALLA}: se rechaza $H_0$.}
  \label{tab:resultados}
  \begin{tabular}{llcccc}
    \toprule
    Generador & Test & Estadístico & Crítico & G.L. & Resultado \\
    \midrule
    \multirow{5}{*}{Cuadrados Medios}
      & Chi-cuadrado ($\chi^2$)    & 10{,}2000  & 16{,}919  & 9  & \textbf{PASA}  \\
      & Kolmogorov-Smirnov ($D$)   & 0{,}027789 & 0{,}043007 & —  & \textbf{PASA}  \\
      & Rachas ($|Z|$)             & 0{,}2531   & 1{,}96    & —  & \textbf{PASA}  \\
      & Autocorrelación ($|Z|$)    & 0{,}7630   & 1{,}96    & —  & \textbf{PASA}  \\
      & Póker ($\chi^2$)           & 50{,}3142  & 12{,}592  & 6  & \textbf{FALLA} \\
    \midrule
    \multirow{5}{*}{GCL (glibc)}
      & Chi-cuadrado ($\chi^2$)    & 6{,}7200   & 16{,}919  & 9  & \textbf{PASA}  \\
      & Kolmogorov-Smirnov ($D$)   & 0{,}020926 & 0{,}043007 & —  & \textbf{PASA}  \\
      & Rachas ($|Z|$)             & 1{,}7718   & 1{,}96    & —  & \textbf{PASA}  \\
      & Autocorrelación ($|Z|$)    & 1{,}2135   & 1{,}96    & —  & \textbf{PASA}  \\
      & Póker ($\chi^2$)           & 4{,}9173   & 12{,}592  & 6  & \textbf{PASA}  \\
    \midrule
    \multirow{5}{*}{Python (MT19937)}
      & Chi-cuadrado ($\chi^2$)    & 8{,}8200   & 16{,}919  & 9  & \textbf{PASA}  \\
      & Kolmogorov-Smirnov ($D$)   & 0{,}031923 & 0{,}043007 & —  & \textbf{PASA}  \\
      & Rachas ($|Z|$)             & 1{,}4554   & 1{,}96    & —  & \textbf{PASA}  \\
      & Autocorrelación ($|Z|$)    & 1{,}2873   & 1{,}96    & —  & \textbf{PASA}  \\
      & Póker ($\chi^2$)           & 6{,}0516   & 12{,}592  & 6  & \textbf{PASA}  \\
    \bottomrule
  \end{tabular}
\end{table}

\subsection{Resumen visual de los tests}

\begin{figure}[H]
  \centering
  \includegraphics[width=0.75\textwidth]{tabla_tests.png}
  \caption{Mapa de calor con el resultado de cada test por generador.
           Verde = PASA, Rojo = FALLA ($\alpha = 0{,}05$).}
  \label{fig:tabla_tests}
\end{figure}

\subsection{Histogramas de uniformidad}

\begin{figure}[H]
  \centering
  \includegraphics[width=\textwidth]{hist_uniformidad.png}
  \caption{Histogramas de frecuencia con $k=10$ intervalos iguales en $[0,1)$.
           La línea discontinua indica la frecuencia esperada ($n/k = 100$).}
  \label{fig:histogramas}
\end{figure}

Los tres generadores muestran distribuciones visualmente similares.  El GCL
y el MT exhiben barras más uniformes, mientras que los Cuadrados Medios
presentan ligeras irregularidades en las frecuencias de los intervalos
extremos, aunque dentro del margen tolerable para el test $\chi^2$.

\subsection{Función de distribución acumulada empírica}

\begin{figure}[H]
  \centering
  \includegraphics[width=0.72\textwidth]{ecdf.png}
  \caption{ECDF de los tres generadores comparada con la distribución
           teórica $F(x)=x$.  Las tres curvas se superponen casi
           perfectamente sobre la diagonal.}
  \label{fig:ecdf}
\end{figure}

La Figura~\ref{fig:ecdf} confirma visualmente el test KS: las tres ECDFs
se adhieren a la teórica sin superar el umbral $D^* = 0{,}043$.

\subsection{Diagrama de dispersión serial}

\begin{figure}[H]
  \centering
  \includegraphics[width=\textwidth]{scatter_serial.png}
  \caption{Diagrama $X_n$ vs $X_{n+1}$ para los primeros 500 pares.
           Una nube uniforme sin estructura indica independencia serial.}
  \label{fig:scatter}
\end{figure}

La nube de puntos del GCL y del MT cubre uniformemente el cuadrado
$[0,1)^2$ sin estructuras apreciables, indicando buena independencia
serial.  Los Cuadrados Medios muestran una distribución también densa,
aunque con una leve concentración diagonal que es consistente con la
autocorrelación no significativa observada.

\subsection{Secuencia temporal}

\begin{figure}[H]
  \centering
  \includegraphics[width=\textwidth]{serie_temporal.png}
  \caption{Primeros 200 valores generados por cada método.  Se aprecia
           el comportamiento errático esperado en todos los generadores.}
  \label{fig:serie}
\end{figure}

\subsection{Autocorrelación para distintos retardos}

\begin{figure}[H]
  \centering
  \includegraphics[width=\textwidth]{autocorrelacion.png}
  \caption{Coeficiente de autocorrelación $r_\ell$ para retardos
           $\ell = 1, \ldots, 50$.  Las líneas discontinuas marcan el
           límite $\pm 1{,}96/\sqrt{n}$ al nivel $\alpha=0{,}05$.}
  \label{fig:autocorr}
\end{figure}

La Figura~\ref{fig:autocorr} muestra que para todos los generadores la
gran mayoría de los coeficientes se encuentran dentro de la banda de
confianza.  En los Cuadrados Medios se aprecian picos esporádicos
marginalmente fuera del límite en algunos retardos, lo que es consistente
con el comportamiento irregular ya detectado por el test del Póker.

\subsection{Estadísticos descriptivos}

\begin{figure}[H]
  \centering
  \includegraphics[width=\textwidth]{estadisticos.png}
  \caption{Media muestral y desviación estándar de cada generador,
           comparadas con los valores teóricos de $U[0,1)$:
           $\mu = 0{,}5$ y $\sigma = 1/\!\sqrt{12} \approx 0{,}2887$.}
  \label{fig:estadisticos}
\end{figure}

Los tres generadores producen medias y desviaciones estándar muy próximas
a los valores teóricos de la distribución uniforme, lo que proporciona
evidencia adicional de la bondad de ajuste.

% ── 5. Análisis comparativo ───────────────────────────────────────────────────
\section{Análisis comparativo}

\subsection{Cuadrados Medios}

El método de los Cuadrados Medios supera los cuatro tests ``clásicos''
(Chi-cuadrado, KS, Rachas y Autocorrelación), pero \textbf{falla el test
del Póker} con un estadístico de $50{,}31$, muy por encima del crítico
$12{,}59$.  Esto indica que la distribución de los dígitos decimales no es
uniforme: hay una sobrerepresentación de ciertas categorías (en particular
``Quintilla'' y ``Trío''), lo que evidencia que los Cuadrados Medios
producen regiones del espacio $[0,1)$ más visitadas que otras cuando se
consideran bloques de cinco valores.

Esta debilidad es bien conocida y motivó históricamente el desarrollo de
generadores más sofisticados.  El método es válido para demostraciones
educativas pero no para simulaciones serias.

\subsection{GCL (glibc/POSIX)}

El GCL con parámetros POSIX supera los cinco tests aplicados.  El test de
Rachas con $|Z|=1{,}77$ (próximo al límite $1{,}96$) sugiere que la
secuencia presenta algo de estructura serial, consistente con la
correlación lineal inherente a los GCLs en dimensiones altas.  No
obstante, para $n=1000$ y los tests empleados, el comportamiento es
estadísticamente aceptable.

Los GCLs tienen período máximo $m = 2^{31} \approx 2{,}1 \times 10^9$,
suficiente para la mayoría de las simulaciones de pequeña y mediana escala.

\subsection{Mersenne Twister (Python)}

El MT supera todos los tests con amplios márgenes.  Es el generador de
mayor calidad de los tres evaluados y el más adecuado para simulaciones de
Monte Carlo de propósito general.  Su único inconveniente para aplicaciones
criptográficas es que su estado interno puede ser reconstruido a partir de
624 valores consecutivos~\cite{matsumoto1998}, pero esto no es relevante
en el contexto de simulación estocástica.

% ── 6. Implementación en Python ───────────────────────────────────────────────
\section{Implementación en Python}

A continuación se muestran los fragmentos más relevantes del código.
El archivo completo se entrega junto con este informe.

\subsection{Generador de Cuadrados Medios}

\begin{lstlisting}[caption={Clase MiddleSquares},label={lst:ms}]
class MiddleSquares:
    def __init__(self, seed: int = 12345, digits: int = 8):
        self.digits  = digits
        self.modulus = 10 ** digits
        self.state   = int(str(seed).zfill(digits)[:digits])

    def next(self) -> float:
        sq   = self.state ** 2
        sq_s = str(sq).zfill(2 * self.digits)
        start = (len(sq_s) - self.digits) // 2
        self.state = int(sq_s[start : start + self.digits])
        return self.state / self.modulus
\end{lstlisting}

\subsection{Generador Congruencial Lineal}

\begin{lstlisting}[caption={Clase GCL},label={lst:gcl}]
class GCL:
    def __init__(self, seed=12345,
                 a=1_103_515_245, c=12_345, m=2**31):
        self.a, self.c, self.m = a, c, m
        self.state = seed % m

    def next(self) -> float:
        self.state = (self.a * self.state + self.c) % self.m
        return self.state / self.m
\end{lstlisting}

\subsection{Test Chi-cuadrado (extracto)}

\begin{lstlisting}[caption={Test de bondad de ajuste Chi-cuadrado},label={lst:chi2}]
def test_chi_cuadrado(data, k=10, alpha=0.05):
    n, Ei = len(data), len(data) / k
    counts = [0] * k
    for x in data:
        counts[min(int(x * k), k - 1)] += 1
    chi2_stat = sum((obs - Ei)**2 / Ei for obs in counts)
    chi2_crit = 16.919  # chi2(9, 0.05)
    return {"estadistico": chi2_stat, "critico": chi2_crit,
            "resultado": "PASA" if chi2_stat < chi2_crit else "FALLA"}
\end{lstlisting}

% ── 7. Conclusiones ───────────────────────────────────────────────────────────
\section{Conclusiones}

Los experimentos realizados permiten extraer las siguientes conclusiones:

\begin{enumerate}

  \item \textbf{Los Cuadrados Medios son insuficientes para simulación.}
        A pesar de pasar los tests de uniformidad global (Chi-cuadrado y
        KS) y los de independencia (Rachas y Autocorrelación), falla el
        test del Póker con un estadístico cuatro veces superior al crítico,
        revelando una distribución no uniforme en la estructura de dígitos.
        Esto lo inhabilita para simulaciones que requieren muestras de
        alta dimensión o bloques de números correlacionados.

  \item \textbf{El GCL con parámetros POSIX es aceptable para simulación
        de uso general.}  Supera los cinco tests aplicados.  Su
        implementación es trivial y su costo computacional es mínimo
        (una multiplicación y una suma módulo $2^{31}$).  La ligera
        tensión en el test de Rachas ($|Z|=1{,}77$) recuerda que los
        GCLs presentan correlaciones en dimensiones altas, limitación
        que los hace inadecuados para métodos de Monte Carlo en espacios
        de alta dimensión.

  \item \textbf{El Mersenne Twister es la opción óptima.}  Supera todos
        los tests con los mejores márgenes y ofrece el período más largo
        ($2^{19937}-1$).  Para cualquier simulación implementada en Python
        se recomienda el uso de \texttt{random.random()} o
        \texttt{numpy.random}, que internamente usan MT19937.

  \item \textbf{Los tests estadísticos son complementarios.}  Ningún test
        aislado es suficiente para certificar la calidad de un generador:
        los Cuadrados Medios pasaron cuatro de los cinco tests.  Se
        recomienda aplicar al menos cuatro pruebas independientes, como
        las aquí presentadas.

  \item \textbf{La semilla determina reproducibilidad.}  Fijar la semilla
        permite reproducir exactamente cualquier simulación, propiedad
        fundamental para la validación y depuración de modelos.

\end{enumerate}

% ── Referencias ───────────────────────────────────────────────────────────────
\begin{thebibliography}{9}

\bibitem{knuth1997}
  Knuth, D.\ E.\ (1997).
  \textit{The Art of Computer Programming, Vol.~2: Seminumerical Algorithms}
  (3rd ed.).
  Addison-Wesley.

\bibitem{metropolis1987}
  Metropolis, N.\ (1987).
  The beginning of the Monte Carlo method.
  \textit{Los Alamos Science}, Special Issue, 125--130.

\bibitem{lehmer1951}
  Lehmer, D.\ H.\ (1951).
  Mathematical methods in large-scale computing units.
  In \textit{Proceedings of the 2nd Symposium on Large-Scale Digital
  Calculating Machinery}, pp.~141--146.
  Harvard University Press.

\bibitem{hull1962}
  Hull, T.\ E.\ \& Dobell, A.\ R.\ (1962).
  Random number generators.
  \textit{SIAM Review}, 4(3), 230--254.

\bibitem{matsumoto1998}
  Matsumoto, M.\ \& Nishimura, T.\ (1998).
  Mersenne Twister: A 623-dimensionally equidistributed uniform
  pseudo-random number generator.
  \textit{ACM Transactions on Modeling and Computer Simulation},
  8(1), 3--30.

\bibitem{massey1951}
  Massey, F.\ J.\ Jr.\ (1951).
  The Kolmogorov-Smirnov test for goodness of fit.
  \textit{Journal of the American Statistical Association},
  46(253), 68--78.

\bibitem{torres2026}
  Torres, J.\ I.\ (2026).
  \textit{TP 2.1 --- Generadores Pseudoaleatorios}.
  Universidad Tecnológica Nacional -- FRRO, Rosario.

\bibitem{tereom2018}
  Organización de Estadística Computacional (2018).
  Números pseudoaleatorios.
  Recuperado de:
  \url{https://tereom.github.io/est-computacional-2018/numeros-pseudoaleatorios.html}

\bibitem{randomorg}
  Haahr, M.\ (2005).
  \textit{Statistical Analysis of RANDOM.ORG}.
  Recuperado de:
  \url{https://www.random.org/analysis/Analysis2005.pdf}

\end{thebibliography}

\end{document}
